\section{The real place: a signed bound for the second case}\label{sec:signed}

In this section we prove the bound of the second case at the rate $U_0$, for configurations with at most
$n/16$ fixed columns. This completes the proof of Theorem~\ref{thm:real}. The section is organized as
follows. The second case contains a determinant $\det\mathsf L$ of size $n$ that the energy method bounds by
$6^n$. We first write it with Pick matrices (\S\ref{ssec:pick}). A Schur complement argument and a contour
bound reduce it to the block of the negative nodes (\S\ref{ssec:schur}--\ref{ssec:T}). Heine's formula turns
that block into a ratio of two integrals of log-gas type (\S\ref{ssec:heine}). We bound the numerator from
above and the denominator from below (\S\ref{ssec:loggas}). Positive nodes close to $1$ need a separate
argument (\S\ref{ssec:spikes}). Combined with OpenAI's energy, the bounds give a concave majorant thanks to an
inequality for signed logarithmic energy (\S\ref{ssec:reflection}). Fixed columns add fixed rows
(\S\ref{ssec:fixedrows}). The hypotheses of the reduction hold on all configurations that matter
(\S\ref{ssec:class}). The section ends with the trials and their verification (\S\ref{ssec:signedcert}).

\subsection{The determinant of the row functions}\label{ssec:pick}
In the second case OpenAI's integrand contains the determinant of the row functions at the points
$x_1,\ldots,x_n\in(-1,1)\setminus\{0\}$ \cite[Prop.~5.3]{OpenAI2026}. With $e=n-1-2g$ it factors as
\begin{equation}\label{eq:detphi}
 \det\varphi(x)=c\,V(x)\prod_a|x_a|^{-(n-1-g)}\det\mathsf L(x),
\end{equation}
where $c$ depends only on $n$ and $g$, and $V(x)=\prod_{a<b}(x_b-x_a)$. To describe $\mathsf L$ let
\[
 B(z)=\prod_a\frac{z-x_a}{1-x_az},\qquad \cK_{ab}=\frac1{1-x_ax_b},
\]
be the Blaschke product of the nodes and the Pick (Cauchy) matrix of the Szeg\H{o} kernels $k_a(z)=1/(1-x_az)$.
Put $w_a=1$ at a negative node and $w_a=-5$ at a positive node, and $D=\operatorname{diag}(w_aB'(x_a)x_a^e)$.
Then $\det\mathsf L=\det(\cK+D)/\det\cK$. In the basis $k_a$ the operator $\mathsf L$ acts on the model space
$\cK_B=H^2\ominus BH^2$, which has dimension $n$. The energy method uses only $\|\mathsf L\|\le6$, which gives
$|\det\mathsf L|\le6^n$. On the configurations near the extremal one, the true size of $\det\mathsf L$ is
smaller by a factor $e^{-cN^2}$. We now capture part of this.

Let $k$ and $q$ be the numbers of negative and positive nodes. For a weight $\omega\in\R$ let
$F(\omega)=\cK+\operatorname{diag}(D_-,\omega D_+)$, where $D_-$ and $D_+$ are the diagonal blocks of
$\operatorname{diag}(B'(x_a)x_a^e)$ at the negative and the positive nodes. So $\det\mathsf L=\det F(-5)/\det\cK$. The case $\omega=1$ is a
Pick matrix.

\begin{lemma}\label{lem:pick}
For every configuration $0\le F(1)\le2\cK$.
\end{lemma}

\begin{proof}
For a real vector $c$ put $\phi(z)=B(z)\sum_ac_a/(z-x_a)$. This is the function of $\cK_B$ with node values
$\phi(x_a)=c_aB'(x_a)$, and $\|\phi\|_{H^2}^2=c^T\cK c$. By the residue theorem on the unit disc
\[
 c^T(F(1)-\cK)c=\sum_a\frac{x_a^e\phi(x_a)^2}{B'(x_a)}=\frac1{2\pi i}\oint_{|z|=1}\frac{z^e\phi(z)^2}{B(z)}\,dz .
\]
On the unit circle $|z^e\phi^2/B|=|\phi|^2$, so the right side lies in $[-c^T\cK c,\,c^T\cK c]$.
\end{proof}

\subsection{A Schur complement bound}\label{ssec:schur}
The determinant $\det F(-5)=\det(F(1)-6\operatorname{diag}(0,D_+))$ is a perturbation of a Pick matrix in the
positive block only. The following inequality bounds it by the negative block.

\begin{theorem}\label{thm:C3}
Let $M=\begin{pmatrix}A&C^*\\C&E\end{pmatrix}$ be a Hermitian block matrix with $0\le M\le b$ and $A>0$. Let $T$
be Hermitian of the size of $E$, with $-\tau\le T\le\tau$, and let $c\ge0$. Then
\[
 \Bigl|\det\begin{pmatrix}A&C^*\\C&E-cT\end{pmatrix}\Bigr|\le(b+c\tau)^{\dim E}\det A .
\]
\end{theorem}

\begin{proof}
The determinant equals $\det A\cdot\det(S-cT)$, where $S=E-CA^{-1}C^*$ is the Schur complement. Since $M\ge0$
one has $S\ge0$, and $S\le E\le b$. So $-c\tau\le S-cT\le b+c\tau$, and every eigenvalue of $S-cT$ has modulus at
most $b+c\tau$.
\end{proof}

We apply this in the metric of $\cK$, that is, after the congruence by a Cholesky factor of $\cK$. By
Lemma~\ref{lem:pick} the hypotheses hold with $b=2$ and $c=6$, provided that $\pm D_+\le\tau\cK$ as quadratic
forms. The result is
\begin{equation}\label{eq:C3}
 |\det\mathsf L|\le(2+6\tau)^q\,\det\mathsf L_{--},\qquad
 \det\mathsf L_{--}=\frac{\det(\cK_{--}+D_-)}{\det\cK_{--}},
\end{equation}
where $\cK_{--}$ is the block of $\cK$ at the negative nodes. The factor $(2+6\tau)^q$ is $e^{O(N)}$. All the
information of order $N^2$ is in $\det\mathsf L_{--}$, which depends on the positive nodes only through the
values $B_+(x_b)$ of their Blaschke product at the negative nodes.

\subsection{A contour bound}\label{ssec:T}
It remains to bound $\tau$. The quadratic form of $D_+$ is the sum of the residues of $z^e\phi^2/B$ at the
positive nodes. These lie in the right half of the disc. We integrate over the boundary of the right
half-disc instead of the circle.

\begin{lemma}\label{lem:T}
Let $M_*=\sup_{-1<y<1}|y|^e/|B(iy)|$. For every real vector $c$
\[
 \Bigl|\sum_{x_a>0}B'(x_a)x_a^ec_a^2\Bigr|\le\Bigl(1+\frac{M_*}2\Bigr)c^T\cK c .
\]
If $\frac en-\frac1n\sum_a\Bigl(\frac{t^2}{t^2+x_a^2}-\frac{x_a^2t^2}{1+x_a^2t^2}\Bigr)>0$ for every
$t\in(0,1]$, then $M_*=1$, and the bound is $\frac32c^T\cK c$.
\end{lemma}

\begin{proof}
With $\phi$ as in the proof of Lemma~\ref{lem:pick}, the sum is the contour integral of $z^e\phi^2/B$ over the
boundary of the right half-disc divided by $2\pi i$. No node lies on this boundary. On the arc the integrand
has modulus $|\phi|^2$, and the arc contributes at most $\|\phi\|^2$. On the diameter $z=iy$ the modulus is at
most $M_*|\phi(iy)|^2$. By the Fej\'er--Riesz inequality on a diameter, $\int_{-1}^1|\phi(iy)|^2dy\le\pi
\|\phi\|^2$. So the diameter contributes at most $\frac12M_*\|\phi\|^2$.

For the second statement write $\log(|y|^e/|B(iy)|)=nh(y)$. Then
$h(y)=\frac en\log|y|+\frac1{2n}\sum_a\log\frac{1+x_a^2y^2}{y^2+x_a^2}$, and $h(1)=0$ because $|B(i)|=1$.
Differentiation gives $th'(t)=\frac en-\frac1n\sum_a(\frac{t^2}{t^2+x_a^2}-\frac{x_a^2t^2}{1+x_a^2t^2})$. If
this is positive on $(0,1]$, then $h<0$ on $(0,1)$. Since $h$ is even, $M_*=1$.
\end{proof}

The formal proof avoids the residue theorem. It subtracts the poles with divided differences, integrates
the remaining holomorphic function over the half-disc with primitives, and proves the Fej\'er--Riesz step
by the same device on the upper half-disc. With $\tau=\frac32$, \eqref{eq:C3} becomes
\begin{equation}\label{eq:eleven}
 |\det\mathsf L|\le11^q\det\mathsf L_{--}.
\end{equation}
The condition on $t$ asks that the node measure is not too concentrated near $0$ at any scale. Its left side
at the extremal measure is at least $0.33$ for every $t$.

\subsection{Heine's formula for the negative block}\label{ssec:heine}
For $s\ge1$ and a positive continuous function $f$ on $[-1,1]$ put
\[
 Z_s(f)=\int_{[0,\pi]^s}\prod_{l<l'}(\cos\theta_l-\cos\theta_{l'})^2\prod_lf(\cos\theta_l)\,d\theta .
\]
By Heine's formula $Z_s(f)$ is $s!$ times the Hankel determinant of the moments of $f(u)(1-u^2)^{-1/2}du$. Let
\[
 \phi(u)=\prod_{x_b<0}(1+x_b^2+2x_bu),\qquad \psi(u)=\prod_{x_p>0}(1+x_p^2+2x_pu)^{-1}.
\]

\begin{theorem}\label{thm:heine}
Let $s=q-g$. If $s\ge1$ and $k\le e+1$, then
\[
 \det\mathsf L_{--}=\prod_{b<c}(1-x_bx_c)\prod_{b,p}(1-x_bx_p)^{-1}\cdot\frac{Z_s(\phi\psi)}{Z_s(\psi)},
\]
where $b,c$ run over the negative nodes and $p$ over the positive nodes.
\end{theorem}

\begin{proof}[Sketch]
One first writes $\det\mathsf L_{--}=\det[x_b^j+B_+(x_b)x_b^{m-j}]_{b,j<k}/V(x_-)$ with $m=e+k-1$, by
Lagrange interpolation at the negative nodes. Since $1+x^2+2xu=x\bigl(x+x^{-1}+2u\bigr)$, the factor $\phi$
is a polynomial in $u$ with roots $-\frac12(x_b+x_b^{-1})$. By Christoffel's formula,
$Z_s(\phi\psi)/Z_s(\psi)$ is a determinant of the monic orthogonal polynomials of the weight
$\psi(u)(1-u^2)^{-1/2}$ at these roots, divided by a Vandermonde determinant. That weight is of
Bernstein--Szeg\H{o} type: $\psi$ is the reciprocal of $|h(e^{i\theta})|^2$ for the polynomial
$h(z)=\prod_p(1+x_pz)$. Its orthogonal polynomials of degree $\ge q$ are explicit. At
$u=-\frac12(x+x^{-1})$ they equal Laurent expressions in $x$ built from $h$. Since $h(z)/z^q\bar h(1/z)$
evaluated at $-x_b$ is $B_+(x_b)$ up to sign, the two determinants agree after the Cauchy-type products are
collected.
\end{proof}

The formal proof works with the angle variables throughout. It needs Heine's formula, Andr\'eief's identity
and the Bernstein--Szeg\H{o} polynomials, but no integration over orthogonal groups. Both integrals in
Theorem~\ref{thm:heine} are of log-gas type. There are $s$ particles $u_l=\cos\theta_l$ on $[-1,1]$ with
logarithmic repulsion, in the external field $\log\phi+\log\psi$ or $\log\psi$. Each field is a sum of
logarithmic potentials $l(x,u)=\log(1+x^2+2xu)$ of the nodes.

\subsection{Bounds for the two integrals}\label{ssec:loggas}
The numerator is bounded from above by the method of OpenAI's energy estimate, now with one population. The
key identity is $\log|u-v|=-\log2-\sum_{k\ge1}\frac2kT_k(u)T_k(v)$ for $u\ne v$ in $[-1,1]$. With the moments
$M_k=\sum_lT_k(u_l)$ the pair sum is $-\binom s2\log2-\sum_k\frac1k(M_k^2-\sum_lT_k(u_l)^2)$. For any sequence
$a$ and any $\nu>0$ one has $-M_k^2\le-2\nu a_kM_k+\nu^2a_k^2$. This completes the square against $a$.

\begin{lemma}\label{lem:LGup}
Let $\nu\ge s$ and let $a$ be a summable sequence. With $\mathsf T(a,u)=\sum_ka_kT_k(u)/k$ and
$\|a\|^2=\sum_ka_k^2/k$,
\[
 \log Z_s(f)\le-s^2\log2+2\nu^2\|a\|^2+s\sup_{[-1,1]}\bigl(\log f-4\nu\mathsf T(a,\cdot)\bigr)+o(\nu^2).
\]
\end{lemma}

The diagonal terms $\sum_lT_k(u_l)^2/k$ diverge when summed over $k$. They are regularized by a damping
$\tau^k$ with $\tau=1-e_N$ as in \cite[Prop.~6.1]{OpenAI2026}, at a cost $o(\nu^2)$. The completed squares
$(M_k-\nu a_k)^2/k$ are dropped in this lemma. They are kept in \S\ref{ssec:class}.

The denominator is bounded from below by Jensen's inequality. Let $\nu'$ be a probability measure on $[0,\pi]$
of the form $\nu'=\bigl(1+2\sum_{k\le K}b_k\cos k\theta\bigr)\frac{d\theta}\pi$, with a positive density. In the
variable $u=\cos\theta$ it is a polynomial multiple of the arcsine law.

\begin{lemma}\label{lem:LGlow}
For every bounded measurable $Q$,
\[
 \log Z_s(e^Q)\ge s(s-1)\Bigl(-\log2-\sum_k\frac{2b_k^2}k\Bigr)+s\int Q(\cos\theta)\,d\nu'(\theta)+s\,h(\nu'),
\]
where $h(\nu')$ is the entropy of the density of $\nu'$ with respect to $d\theta$.
\end{lemma}

\begin{proof}
Apply Jensen's inequality to the product density $\prod_l\frac{d\nu'}{d\theta}(\theta_l)$. The logarithmic
energy of $\nu'$ in the variable $u$ is $-\log2-\sum_k2b_k^2/k$, by the expansion of $\log|u-v|$ above.
\end{proof}

The fields of the nodes integrate against $\nu'$ in closed form. For $|x|<1$ one has
$\int l(x,\cos\theta)\,d\nu'=-2\sum_k(-x)^kb_k/k$, a polynomial in $x$. The arcsine part contributes $0$,
because the logarithmic potential of the arcsine law is constant on $[-1,1]$.

\subsection{Positive nodes near $1$}\label{ssec:spikes}
A positive node $x_p=1-\epsilon$ puts into $\psi$ a factor $(1+x_p^2+2x_pu)^{-1}$ of size $\epsilon^{-2}$ near
$u=-1$. Lemma~\ref{lem:LGup} charges the supremum of the field to all $s$ particles. A single such node
would therefore cost about $N\log(1/\epsilon)$, while the true change of $\det\mathsf L_{--}$ is bounded. The
following reduction removes such nodes from the log gas.

\begin{theorem}\label{thm:spike}
Let $\phi$ be positive and non-increasing on $[-1,1]$, and let $\psi=\psi_r\prod_{p\in P}k_p$ with $\psi_r>0$,
$k_p(u)=(1+2y_pu+y_p^2)^{-1}$ and $y_p\in(0,1)$. If $|P|=2m$ or $|P|=2m-1$, then
\[
 \log\frac{Z_s(\phi\psi)}{Z_s(\psi)}\le m\log\phi(-1)+\log\frac{Z_{s-m}(\phi\psi_r)}{Z_{s-m}(\psi_r)} .
\]
\end{theorem}

\begin{proof}[Sketch]
For an odd number of factors add one more factor, at a point to be chosen. Write both integrals as Hankel
determinants. Moving all $y_p$ to a common value changes the ratio monotonically. This uses a one-point
comparison for the weight $(1+2yu+y^2)^{-1}$, whose Chebyshev moments form the geometric sequence $(-y)^k$.
For a common value $y$ the $2m$ factors behave like $m$ particles pinned near $u=-1$. Fischer's inequality
bounds the numerator by $\phi(-1)^m$ times the determinant of the remaining block. As $y\to1$ the Schur
complement of the pinned block in the denominator converges to the Hankel matrix of the remaining
$s-m$ particles.
\end{proof}

We apply the theorem to the positive nodes with $1-x_p<\theta_0=10^{-3}$. Since $\phi(-1)=\prod_b(1-x_b)^2\le
4^k$, each such pair costs at most $O(N)$. The energy of OpenAI's estimate charges each such node
$\beta_hn\log(1/\theta_0)$, which is larger. So the number of such nodes on the configurations that matter
is bounded by a constant times $N$, with a cost absorbed in the margin.

\subsection{Concavity: an inequality for signed logarithmic energy}\label{ssec:reflection}
We now combine OpenAI's energy with \eqref{eq:eleven}, Theorem~\ref{thm:heine} and
Lemmas~\ref{lem:LGup} and~\ref{lem:LGlow}. The total is a functional of the node measure $\mu$, split into its
negative part $\mu_-$ and positive part $\mu_+$, and of the column measure. Its quadratic part is not obviously
negative. Theorem~\ref{thm:heine} contributes the Cauchy-type term $-\iint\log(1-xy)\,d\mu_-(x)d\mu_+(y)$,
and Lemma~\ref{lem:LGup} applied with a trial that depends on $\mu$ contributes a convex term. Write
\[
 I(\delta)=-\iint\log|x-y|\,d\delta(x)d\delta(y),\qquad L(\tau)=-\iint\log(1-xy)\,d\tau(x)d\tau(y)
\]
for real measures on $[-1,1]$ and $[0,1]$. Both are sums of squares of moments: $L(\tau)=\sum_k\frac1k\bigl(\int
x^kd\tau\bigr)^2$. The quadratic part of the combined functional is $-\frac12I-\frac12\sum_k(B_k-2C_k)^2/k+\frac12
L(\mu_+)$, where $B_k$ and $C_k$ are the power moments of the nodes and the columns. The term $\frac12L(\mu_+)$
is absorbed by the following theorem.

\begin{theorem}\label{thm:reflection}
Let $\delta=\delta_-+\delta_+$ be a real measure of finite logarithmic energy and total mass $0$, with
$\delta_-$ on $[-1,0]$ and $\delta_+$ on $[0,1]$. Let $R$ be the reflection $x\mapsto-x$. Then
\[
 I(\delta)\ge L(\delta_+)+L(R\delta_-)\ge L(\delta_+).
\]
\end{theorem}

\begin{proof}
Put $p=\delta_+$ and $r=R\delta_-$, both on $[0,1]$, and let $K_1(a,b)=-\log|a-b|$ and $K_2(a,b)=-\log(a+b)$.
Since $|x-y|=a+b$ for nodes of opposite signs, $I(\delta)=K_1(p)+K_1(r)+2K_2(p,r)$, which equals
$\frac12(K_1+K_2)(p+r)+\frac12(K_1-K_2)(p-r)$.

First, $p+r$ has mass $0$. One has $K_1=G_{\mathbb D}+L$, where $G_{\mathbb D}(a,b)=\log|(1-ab)/(a-b)|$ is the
Green function of the unit disc, whose energy is non-negative. By Frullani's integral
$-\log(a+b)=\int_0^\infty(e^{-(a+b)t}-e^{-t})\frac{dt}t$, so $K_2(\sigma)=\int_0^\infty(\int e^{-at}d\sigma(a))^2
\frac{dt}t\ge0$ for $\sigma$ of mass $0$. Hence $(K_1+K_2)(p+r)\ge L(p+r)$.

Second, $H=K_1-K_2=\log\frac{a+b}{|a-b|}$ satisfies $H(\sigma)\ge L(\sigma)$ for every real $\sigma$ on $[0,1]$.
With $a=e^{-x}$ and $b=e^{-y}$ one has $H=\log\coth\frac{|x-y|}2=\sum_{j\text{ odd}}\frac2je^{-j|x-y|}$. For
$x,y\ge0$, $e^{-j|x-y|}=e^{-j(x+y)}+2j\int_0^{\min(x,y)}e^{-j(x+y-2v)}dv$. Put $F_j(v)=\int_{x\ge v}e^{-j(x-v)}
d\sigma$. Then $H(\sigma)=\sum_{j\text{ odd}}\bigl(\frac2jF_j(0)^2+4\int_0^\infty F_j(v)^2dv\bigr)$. In
$L(\sigma)=\sum_k\frac1k(\int e^{-kx}d\sigma)^2$ the terms with odd $k$ are $F_k(0)^2/k$. For even $k=j+1$,
integration by parts gives $\int e^{-(j+1)x}d\sigma=F_j(0)-\int_0^\infty e^{-(j+1)v}F_j(v)\,dv$. The
Cauchy--Schwarz inequality with weights $1+\frac1j$ and $1+j$ gives $(\int e^{-(j+1)x}d\sigma)^2\le
(1+\frac1j)F_j(0)^2+\frac12\int F_j^2$. Summing, $L(\sigma)\le H(\sigma)$.

Together, $I(\delta)\ge\frac12L(p+r)+\frac12L(p-r)=L(p)+L(r)$.
\end{proof}

By Theorem~\ref{thm:reflection} the quadratic part is at most
$-\frac14I-\frac12\sum_k(B_k-2C_k)^2/k\le0$ on differences of admissible measures. So the combined functional
is concave. Its maximum is bounded by the value at a trial plus the supremum of a linear functional, that
is, by a constant plus the suprema of fields of one variable. The formal proof uses a discrete form of
Theorem~\ref{thm:reflection} for weighted point configurations against a reference measure with bounded
density. The fixed rows of the next subsection change the masses of the two species. The missing mass is
compensated by phantom particles distributed by the arcsine law, whose potential is constant on $[-1,1]$.

\subsection{Fixed rows}\label{ssec:fixedrows}
With $m$ fixed columns the determinant of the row functions is evaluated at the free points and at the fixed
points. The fixed points enter $\det\mathsf L$ as nodes, but they are not integrated over. Lemma~\ref{lem:T}
then needs a change. Its constant $\frac32$ is false when the sum runs over the free positive nodes only, while
the Blaschke product and the Gram form contain the fixed nodes as well. For two adjacent positive nodes $0.5$
and $0.55$ with $e=2$ the best constant is $3.625$. The correct statement keeps the fixed positive residues as a
separate term $W=\sum x_a^e/((1-x_a^2)|B'(x_a)|)$ over the fixed positive nodes. Each of them is bounded by the
reproducing estimate $(c_aB'(x_a))^2(1-x_a^2)\le c^T\cK c$.

Instead of absorbing $W$ into $\tau$, we expand. Write $\det(X+\operatorname{diag}\delta)$ as a sum of principal
minors over subsets $S$ of the fixed positive nodes. This gives
\[
 |\det\mathsf L|\le11^q\sum_S\Bigl(\frac6{11}\Bigr)^{|S|}\det\mathsf L_{--}^{(S)}\,W_S,
\]
where in the term for $S$ the nodes of $S$ interact with the other nodes through $\log(1-xy)$ instead of
$\log|x-y|$. The energy then has three fixed populations: fixed negative nodes, fixed positive nodes, and the
nodes of $S$. Their fields are bounded on their domains like the residual field of Section~\ref{sec:real}.
The fixed nodes are images of quadrature nodes. So the mass of their measure $\kappa$ near $1$ is limited
by the spacing of the quadrature nodes, namely $\kappa([x,1])\le C_1(1-x)$ for $x\ge x_c=0.999$. For such a measure of mass $\alpha$
and a field $\psi$ with $\psi\le\lambda$ on $[0,x_c]$ and $\psi(x)\le g_0+r\log\frac1{1-x^2}$ on $[x_c,1)$,
\[
 \int\psi\,d\kappa\le\alpha\lambda+C_1\int_0^{1-x_c}\bigl(g_0+r\log(1/u)-\lambda\bigr)_+\,du .
\]
The integral has a closed form. This bounds the fixed rows without a finite optimization over their positions.

\subsection{The configurations that matter}\label{ssec:class}
The reduction above needs hypotheses on the configuration. Lemma~\ref{lem:T} needs the condition on $t$. The
number of positive nodes within $\theta_0$ of $1$ must be bounded by a constant times $N$. The fraction of
negative nodes must lie in a fixed window. These hypotheses hold on every configuration whose energy, in
OpenAI's estimate, is within the gain of the maximum. We use the following dichotomy. Either OpenAI's bound at
the rate $U_0$ holds for the configuration, and nothing is to be proved. Or the configuration exceeds it. Then
its energy is within $U_0-U$ of the maximum of OpenAI's bound.

For such a configuration, keep the completed squares and the non-negative differences between the fields and
their suprema in the proof of OpenAI's estimate. Their sum is at most the deficit, a constant times $N^2(U_0-U)$.
This bounds the fraction of nodes on which a field is far below its supremum, by Chebyshev's inequality. It
also bounds every linear statistic $\sum_kF_k(M_k-\nu a_k)/k$ of the moments, by the Cauchy--Schwarz
inequality. The condition of Lemma~\ref{lem:T} is a family of such linear statistics in $t$. It is checked on a
grid of $24$ values of $t$, with a polynomial lower bound between the grid points. The count of nodes near $1$
follows from the field of the positive nodes near $x=1$, which lies below its supremum by a fixed amount. The
window for the fraction of negative nodes follows from polynomial minorants of the indicator functions of the
two signs.

\subsection{The trials and their verification}\label{ssec:signedcert}
The concave majorant of \S\ref{ssec:reflection} is bounded at a trial. A trial consists of a reference measure
for the free nodes, with a piecewise polynomial density, trial sequences for the rows, the columns and the
numerator of Theorem~\ref{thm:heine}, and the measure $\nu'$ of Lemma~\ref{lem:LGlow}. The trial sequences
have $120$ Chebyshev coefficients and geometric tails with complex poles, as in \eqref{eq:trial}. The bound
then reads
\begin{multline*}
 \text{constant}+\alpha\sup\text{(fixed-row fields)}+\sup\text{(free-node fields)}\\
 +\sup\text{(column field)}\le-U_0+\tfrac12\log2+c_2\alpha .
\end{multline*}
The constant is the value at the trial. It consists of logarithmic energies and moments of the reference
measure, the norms of the trial sequences and the entropy of $\nu'$. All of these have closed forms.

The fraction $\alpha$ of fixed columns varies in $[0,\frac1{16}]$. We use trials at $\alpha=0$, $\frac1{200}$ and
$\frac1{16}$. On each of the two intervals between them we interpolate the trials linearly in $\alpha$, in all
the data weighted by mass. Every field is then exactly affine in the interpolation parameter. So its supremum
on the interval is bounded by the chord of its suprema at the two ends. The constant is a quadratic polynomial
in the parameter. Its curvature is an explicit combination of the four cross energies of the two end trials.
The spike constant of Theorem~\ref{thm:spike} interpolates in the same way. At each end trial we therefore
bound the suprema of the following fields: the fields of the negative and the positive free nodes, the
column field, and the fields of the fixed negative, fixed positive and expanded fixed nodes. Each is a
polynomial plus finitely many terms $P(x)\log|x-c|$ and $\operatorname{Re}\rho\log(x-z)$ with complex $z$, plus
a constant.

The suprema are verified on covers by cells, as in \S\ref{ssec:cells}, with the endpoints of the domains
treated by Laurent expansions and logarithmic end windows. The polynomial part of a field has degree about
$120$. On a block of cells we shift its Taylor expansion exactly to order $K$ and bound the remaining
coefficients by one table shared by the block. This keeps the integers small. The three end trials need
\NcellsSigned{} cells in total. The smallest margin of the budget inequality is $0.037$ per unit of $N^2$.

This completes the proof of the bound of the second case at the rate $U_0$ for $\alpha\le\frac1{16}$. With
\S\ref{ssec:cases} it proves Theorem~\ref{thm:real}.
